\section{Results, Discussion, and Future Work}
\label{sec:results}

This chapter summarises what is verified, what is not yet
verified but planned, what the firmware deliberately does not
do, and the staged path between the present implementation and
a flight-ready EPS.

\subsection{What is verified}

The following items are verified to be working on the real PCU
testing board V4.1 at the time of writing:

\begin{itemize}
\item The board boots, the 48\,MHz clock initialises, the
status LED blinks, and the watchdog is fed by the main loop.
\item The CHIPS link on J3 exchanges frames bi-directionally
with the ESP32 bridge at 115200\,baud. The bus scan, the
\texttt{telemetry}, the \texttt{state}, the \texttt{adc}, the
injected-sensor commands, and the PWM safety commands all
return valid responses.
\item The twelve state-transition scenarios pass against the
pure-dispatcher firmware (12/12 green).
\item The MPPT demo runs in closed loop against the ESP32-side
physics model. Convergence is visually clear on the web page;
the duty cycle settles within tens of iterations.
\item The complementary PWM at 300\,kHz on PA12/PA13 has been
scoped at one edge, with a dead-time gap of approximately
100\,ns. The PWM safety gate disarms on boot, on mode change,
and on any injected fault.
\item Manual mode toggles each of the four user-facing outputs
on the dev board, with the snapshot mirror behaviour matching
the requested values even on outputs that are not physically
wired on the dev board.
\end{itemize}

\subsection{What the bench programme will verify}

The nine-test bench programme described in
Chapter~\ref{sec:hil}.\ref{sec:bench-test} produces the
remaining evidence. Once executed, the following statements
will be supported:

\begin{itemize}
\item The PWM dead-time is consistent across the duty range and
on both edges. (Test A.)
\item Every ADC sensor channel reads correctly against a known
applied voltage, with calibration coefficients quoted. (Test B.)
\item Each eFuse can be enabled and disabled by firmware
commands; the post-eFuse rail and the PGOOD telemetry track the
commanded state. (Test C.)
\item The state machine reaches the expected mode for each
sensor combination when fed by real ADC instead of injected
values. (Test D.)
\item The MPPT algorithm running on the real chip converges to
within a small percentage of the hand-calculated maximum power
point of a fake-solar source (bench supply with a series
resistor). (Test E --- the headline figure of the report.)
\item The buck stage follows the ideal step-down relationship
across the operating duty range. (Test F.)
\item The two INA226 chips respond on the I\textsuperscript{2}C
bus at their schematic-derived addresses and produce readings
that agree with a multimeter. (Test G.)
\item The CHIPS link tolerates sustained streaming traffic
without data loss and rejects corrupted frames as required by the
protocol. (Test H.)
\item The watchdog fires on a deliberate stall and the board
recovers cleanly. (Test I.)
\end{itemize}

A blank table that the user will fill in as each test is
executed is included in the appendix.

\subsection{What the simulator deliberately does not prove}
\label{sec:not-proven}

The state machine running today is a single short pure
function that maps sensor values to a PCU mode. It is the
simplest implementation that satisfies the steady-state
specification, and it deliberately omits the per-mode
machinery that flight-grade firmware will need. The
\texttt{state-transitions.md} file inside the repository
enumerates these omissions; the most important are:

\begin{itemize}
\item Timer-based transitions. The flight FSM moves from
\texttt{MPPT\_CHARGE} to \texttt{CV\_FLOAT} only after the
battery has been at full voltage for at least \texttt{t1}
consecutive cycles. The current dispatcher detects
``battery full'' instantaneously.

\item The CV\_FLOAT \texttt{TEMP\_MPPT} sub-state. A transient
large load that briefly drains the battery while CV\_FLOAT is
active should trigger a temporary return to the MPPT
algorithm. The current dispatcher does not implement this.

\item Per-mode duty algorithms. Each mode in flight uses its
own control law for the buck duty cycle (incremental
conductance in \texttt{MPPT\_CHARGE}, bang-bang voltage
regulation in \texttt{CV\_FLOAT}, current-clamped operation in
\texttt{SA\_LOAD\_FOLLOW}, minimum-duty in
\texttt{BATTERY\_DISCHARGE}). The simulator uses placeholder
duty values for everything but \texttt{MPPT\_CHARGE}.

\item Hysteresis and noise robustness. Real sensor readings
jitter. The flight FSM debounces threshold crossings; the
simulator receives clean injected values.

\item Load shedding inside \texttt{BATTERY\_DISCHARGE}. When
the discharge current exceeds the safe limit the flight FSM is
supposed to shed loads one at a time in priority order. The
simulator does not implement this.

\item Fault-flag interpretation. The ESP32 bridge can inject a
fault bitmask; the current dispatcher does not yet decode any
bit of it.

\item Tuned threshold values. Every numeric threshold (battery
maximum voltage, battery minimum voltage, temperature limits,
heartbeat timeout) is a placeholder. Real values depend on the
chosen battery cell, the thermal model, and the orbit profile,
and will be set when the relevant hardware decisions are made.
\end{itemize}

These omissions are not bugs: they are explicit scope choices
that keep the simulator-grade firmware small enough to be
inspected and verified end-to-end during the project's
duration.

\subsection{Path to flight code}

The work remaining to grow the present implementation into a
flight-grade EPS is itemised in the
\texttt{what\_the\_simulator\_does\_and\_what\_real\_flight\_software\_still\_needs.md}
file in the repository. The eleven steps are summarised here:

\begin{enumerate}
\item Add a control-loop iteration counter to the firmware's
shared state. This is the foundation for every timing-based
behaviour.
\item Implement the \texttt{t1} timeout for the
\texttt{MPPT\_CHARGE} $\rightarrow$ \texttt{CV\_FLOAT}
transition.
\item Implement the \texttt{t2} debounce for the return
transition.
\item Add the \texttt{TEMP\_MPPT} sub-state inside
\texttt{CV\_FLOAT}.
\item Replace the placeholder duty cycle in
\texttt{MPPT\_CHARGE} with the real Incremental Conductance
algorithm. (Already implemented in pure C; only the
integration is pending.)
\item Implement bang-bang voltage regulation inside
\texttt{CV\_FLOAT}.
\item Implement the current-clamp duty algorithm inside
\texttt{SA\_LOAD\_FOLLOW}.
\item Implement per-iteration load shedding inside
\texttt{BATTERY\_DISCHARGE}.
\item Add low-pass filters and hysteresis bands for every
sensor channel.
\item Define and interpret the fault-flag bits.
\item Replace every placeholder numeric threshold with the
mission-tuned values once the corresponding hardware is
finalised.
\end{enumerate}

Each of these can be implemented and verified by adding a new
scenario to the existing twelve-scenario web page; the test
infrastructure does not need to change.

\subsection{Open issues}

A small set of mission-level open questions affect the
firmware but cannot be answered inside the firmware itself:

\begin{itemize}
\item The battery cell has not been chosen, so the values of
$V_{\text{bat,max}}$, $V_{\text{bat,full}}$,
$V_{\text{bat,min}}$, $V_{\text{bat,critical}}$, the maximum
charge and discharge currents, and the temperature limits are
all placeholders.
\item The SPI-attached battery-temperature sensor is named in
the mission document but the exact part is not specified, so
its driver is currently a stub.
\item The mission document and the team's communication
choices disagree on whether the OBC--EPS link is
I\textsuperscript{2}C or UART; the current firmware uses UART
per the most recent team decision.
\item The PCU V4.1 PWM routing requires the dual-channel TCC0
dead-time workaround documented in
\texttt{pcb\_design\_error\_for\_pwm.md}; this is implemented
in firmware but will need to be revisited when the next PCB
revision is laid out.
\item The PCB design contains a small documentation
discrepancy: the EPC2152 reference designator on the
silkscreen is \texttt{IC6}, whereas an older revision of the
mainboard pinout document refers to it as \texttt{U1}. The
firmware itself is unaffected; the documentation has been
updated.
\end{itemize}

\subsection{Conclusion}

The work presented in this report delivers a firmware
architecture for the CHESS EPS that is testable end-to-end on
real hardware, decoupled by design between pure logic and
hardware access, and structured so that the path from the
present simulator-grade implementation to a flight-grade EPS
is a sequence of well-scoped additions rather than a rewrite.
The hardware-in-the-loop test harness, the twelve state
scenarios, the MPPT convergence demonstration, and the
nine-test bench programme together provide a level of
verification coverage that is unusual for a one-semester
student project, and that gives the next student to inherit
the work a clean starting point.

The next student should source the power resistors required
for the bench programme, execute the nine tests, and then
begin the eleven-step climb to flight code. Each step is a
single isolated change with a single new test scenario, and
the architecture is ready to absorb them.
